\section{Conclusions}
\label{sec:conclusion}

This paper evaluated gravity and magnetic anomaly aiding of a consumer-grade MEMS INS under degraded GNSS with three estimators on 27 drives, pairing every aided run with the same filter's unaided run on an identical measurement stream. It corrects the authors' earlier conference results, which arose from unit errors in the anomaly observation and from comparisons against differing or diverged baselines.

With corrected observations, the gravity and magnetic readings that a vehicle-mounted smartphone already provides carry no map information and provide no benefit to either Kalman filter. The cause is the readings themselves: the gravity norm is pinned to a device constant and the magnetometer is dominated by the vehicle's field, which gives signal-to-noise ratios of 0.13 and 0.01 against the maps. Replacing only those readings with a simulated \$25 RM3100 magnetometer makes magnetic aiding reduce the median error of both the EKF and the UKF by about 10\%, on 25 and 26 of 27 drives, while a simulated \$77 ADXL355 used as a gravimeter does not help. The size of the magnetic gain is consistent with a single-update heuristic, the sensor noise divided by the map gradient. The deciding variable is the quality of the anomaly measurement relative to the map signal, not the estimator and not the inertial sensor; by the same heuristic, an indicative target on these roads and maps is a noise of about \boundreqgrav~mGal or \boundreqmag~nT.

The Rao-Blackwellized particle filter, which has the lowest median error of the three with full GNSS, loses track with one fix per minute with or without aiding, in a manner consistent with particle deprivation, and its results are far less repeatable than those of the Kalman filters. In the regime studied, where the position uncertainty is much smaller than the correlation lengths of the maps, the Kalman filters are the estimators that worked; whether the particle filter can be made to work under sparse GNSS is open.

Four directions follow. A field trial with a dedicated magnetometer mounted away from, or compensated for, the vehicle's field would test whether the simulated gain survives installation. The Kalman filters' measurement Jacobian should include the reference-field and motion terms, and where the anomaly errors are correlated in time, as the smartphone's vehicle field is, that correlation should be modeled or the update interval lengthened to match it. The particle filter needs a formulation that survives sparse GNSS, by updating position linearly or proposing particles near each fix, before its multimodal advantage can be tested in long outages. And the error-state Kalman filter that \texttt{strapdown-rs} provides as its default estimator should be given an anomaly path and evaluated alongside the filters compared here.
